\section{Introduction}
\label{sec:intro}

The chest radiograph is the most frequently acquired imaging examination and the first-line test for many thoracic conditions, including pneumonia, pleural effusion, pulmonary nodules and tuberculosis (TB). Reading radiographs requires trained readers, who are scarce in many of the settings where TB and pneumonia are most common. Automated classification of chest radiographs has therefore attracted sustained attention, from early TB screening systems built on small curated sets~\cite{jaeger2014two,candemir2014lung} to deep convolutional networks trained on tens of thousands of hospital images~\cite{wang2017chestxray8,rajpurkar2017chexnet,irvin2019chexpert,johnson2019mimic}. The COVID-19 pandemic added a wave of models that classify radiographs into COVID-19 and other categories~\cite{wang2020covidnet,apostolopoulos2020covid,chowdhury2020can}.

Many published multi-class chest radiograph classifiers are trained and tested on data pooled from several public collections, because no single public collection covers all the classes of interest. Pooling makes three failure modes likely. First, \emph{patient leakage}: several images of the same patient can fall into both the training and the test split when splits are drawn per image, or when the same image reaches a project through two redistributed archives. Second, \emph{proxy labels}: a class can be filled with images whose source label means something else, such as ``abnormal, not TB'' being used as ``pneumonia''. Third, \emph{source confounding}: when a class comes mainly or only from one collection, the network can learn the collection's acquisition style, text markers or patient age instead of the disease, and internal test images from the same collection reward that shortcut~\cite{zech2018variable,geirhos2020shortcut,degrave2021ai,maguolo2020critic}. Systematic reviews of COVID-19 imaging models found these problems to be widespread and concluded that very few published models were fit for clinical use~\cite{wynants2020prediction,roberts2021common}. Reporting guidelines for medical imaging artificial intelligence now ask for explicit descriptions of data sources, partitioning and external evaluation~\cite{mongan2020claim,collins2024tripodai}.

This study began as an audit of our own earlier pipeline (here called V5). V5 pooled eight public sources into a six-class dataset and reported a high internal test accuracy (\RVFiveAcc\% on $N=\RVFiveN$ test images). The audit, summarised in Section~\ref{sec:disc-audit}, found that 154 test images shared a patient with training images, that part of the Pneumonia and Tuberculosis classes consisted of proxy or latent-disease labels, and that on an external hospital cohort the model called none of the \RVFiveExtSpecN{} normal radiographs Normal. The high internal number was therefore not evidence of a useful classifier. We rebuilt the dataset from official sources under rules fixed before training (V6), retrained the classifier, and evaluated it with the external cohort and per-source analyses that V5 lacked.

The question this paper addresses is not ``how accurate can a six-class chest radiograph classifier be'', but ``how much of the internal accuracy of a multi-source chest radiograph classifier survives a leakage-controlled split, per-source analysis and an unseen hospital''. We make no claim of clinical validity.

\paragraph{Contributions}
\begin{enumerate}
\item \textbf{An auditable, leakage-controlled rebuild of a multi-source six-class chest radiograph benchmark from official sources.} Every file is logged with its download URL, size and SHA-256 digest; labelling rules, exclusions and per-source caps are versioned in a configuration file and fixed before training; splits are grouped by patient where identity is known and by merged exact and perceptual-hash near-duplicates otherwise; the build report records zero groups and zero identical files spanning splits, and no external-cohort file in the development data (Section~\ref{sec:materials}). We state explicitly where patient identity is unknown and separation is therefore not proven.
\item \textbf{A reproducible end-to-end pipeline} from acquisition through grouped splitting, training, once-only test and external evaluation, to the metrics served by a web application, with checksums and split fingerprints carried through every stage (Sections~\ref{sec:methods} and~\ref{sec:design}).
\item \textbf{Evidence that internal metrics overstate generalisation.} We report internal results together with per-source accuracy, the predicted-class breakdown on an external TB cohort, and external sensitivity always with specificity (Section~\ref{sec:results}). Although its internal macro ROC-AUC was \RTestMacroAUC, the selected model called only \RExtSpecK/\RExtSpecN{} normal radiographs from the unseen hospital Normal and named \RExtTBK/\RExtTBN{} TB radiographs Tuberculosis, its test accuracy ranged from \RSrcAccCohen\% to \RSrcAccCovidRad\% across sources, and a linear probe recovered the source of an image from frozen features with \RProbeImgBalAcc\% balanced accuracy (chance \RProbeImgChance\%), also within single classes (Section~\ref{sec:res-probe}).
\item \textbf{An audit of an earlier pipeline} showing how patient leakage, proxy labels and the absence of external specificity reporting inflated a recorded result, and a deployment-level finding that a heuristic input gate rejected most genuine radiographs (Section~\ref{sec:discussion}).
\end{enumerate}

We do not propose a new network architecture. The classifier is a standard ImageNet-pretrained DenseNet-121~\cite{huang2017densely,deng2009imagenet} so that the effect of data construction and evaluation protocol is not confounded with architectural novelty. The main results come from a single training run per dataset build (seed \RSeed), on a laptop CPU, under a time budget. Two further seeds and a preprocessing ablation of the selected configuration, fixed in advance, are reported in Sections~\ref{sec:res-seeds} and~\ref{sec:res-ablation}.

The rest of the paper is organised as follows. Section~\ref{sec:related} reviews related work. Section~\ref{sec:materials} describes the sources and the dataset construction. Section~\ref{sec:methods} gives the preprocessing, model, training and metrics. Section~\ref{sec:design} describes the experimental design, including the two dataset builds and the rule used to choose between them. Section~\ref{sec:results} reports results, Section~\ref{sec:discussion} discusses them together with the audit and the limitations, and Section~\ref{sec:conclusion} concludes.
